\documentclass[a4paper,12pt]{article}
\usepackage[applemac]{inputenc}
\usepackage[OT1,T1]{fontenc}

\setlength{\textwidth}{160mm}
\setlength{\textheight}{230mm}
\setlength{\oddsidemargin}{-5mm}
\setlength{\topmargin}{-10mm}
% \usepackage{times}
\usepackage{setspace}
\usepackage{lscape}
\usepackage{amsmath,amssymb}
\usepackage{fancyvrb}
\usepackage{multirow}
\usepackage{float}
\usepackage{color}
\usepackage{listings}
\usepackage{fancyhdr}

\usepackage{fancyvrb,moreverb,boxedminipage}

\textwidth=6in
\textheight=8in
\topmargin=0.0in
\oddsidemargin=0in 
\baselineskip=18pt
\headheight=2cm
\setcounter{MaxMatrixCols}{10}

\newcommand{\ns}{\!\!\!}
\newcommand{\e}{\varepsilon}
\newcommand{\be}{\beta}
\newcommand{\s}{\sigma}
\newcommand{\vv}{\vspace{.5cm}}
\DeclareMathOperator{\plim}{plim}
\DefineVerbatimEnvironment{code}{Verbatim}{fontsize=\small}
\DefineVerbatimEnvironment{example}{Verbatim}{fontsize=\small}

\def\by{\mathbf{y}}
\def\bx{\mathbf{x}}
\def\ba{\mathbf{a}}
\def\bb{\mathbf{b}}
\def\be{\mathbf{e}}
\def\bu{\mathbf{u}}
\def\bv{\mathbf{v}}
\def\bz{\mathbf{z}}

%===========
%Greek letter vectors
%===========
\def\bbeta{\mbox{\boldmath $\beta$}}
\def\bgamma{\mbox{\boldmath $\gamma$}}
\def\bvareps{\mbox{\boldmath $\varepsilon$}}
\def\bleta{\mbox{\boldmath $\eta$}}
\def\bmu{\mbox{\boldmath $\mu$}}
\def\btheta{\mbox{\boldmath $\theta$}}
\def\bsigma{\mbox{\boldmath $\sigma$}}
\def\blambgam{\mbox{\boldmath $\lambda_{\gamma}$}}
\def\blambbet{\mbox{\boldmath $\lambda_{\beta}$}}
\def\blambda{\mbox{\boldmath $\lambda$}}

%===========
%Matrices
%===========
\def\bX{\mathbf{X}}
\def\bD{\mathbf{D}(\bsigma)}
\def\bV{\mathbf{V}(\bsigma)}
\def\bVinv{\mathbf{V}^{-1}(\bsigma)}
\def\bR{\mathbf{R}(\bsigma)}
\def\bA{\mathbf{A}}
\def\bB{\mathbf{B}}
\def\bZ{\mathbf{Z}}
\def\bIn{\mathbf{I_{N}}}
\def\bIm{\mathbf{I_{m}}}
\def\bIni{\mathbf{I_{n}}}
\def\SSW{\mathrm{SSW}}
\def\SSB{\mathrm{SSB}}
\def\Normal{\mathcal{N}\!}

%===========
%Parentheses, etc.
%===========
\def\op{\left(}
\def\cp{\right)}
\def\oc{\left[}
\def\cc{\right]}
\def\ok{\left\{}
\def\ck{\right\}}

%===========
%Statistical functions and real numbers.
%===========
\def\exE{\mathbb{E}}
\def\Real{\mathbb{R}}
\def\Var{\mathbb{V}ar}
\def\Proba{\mathbb{P}}


\renewcommand{\labelenumi}{\bfseries\alph{enumi})}
\newcommand{\splus}[3][0]{\texttt{>~#2}\ \ \ \ \textit{#3}\\[#1mm]}

\newenvironment{exo}
{\begin{center}\setlength{\textwidth}{140mm} \underline{Exercise:} \begin{minipage}[t]{120mm}}
{\end{minipage}\setlength{\textwidth}{160mm}\end{center}}


\usepackage{verbatim}

\lstset{ %
  language=R, 
  basicstyle=\ttfamily,
  backgroundcolor=\color{lightgray},   % choose the background color; you must add \usepackage{color} or \usepackage{xcolor}
  xleftmargin= 10 pt,
  xrightmargin= 10 pt,
  escapeinside={\%*}{*)},          % if you want to add LaTeX within your code
  frame=none,                    % adds a frame around the code
  numbers=none,                    % where to put the line-numbers; possible values are (none, left, right)
  numbersep=5pt,                   % how far the line-numbers are from the code
  numberstyle=\tiny\color{mygray}, % the style that is used for the line-numbers
  showspaces=false,                % show spaces everywhere adding particular underscores; it overrides 'showstringspaces'
}

\pagestyle{fancy}

\definecolor{mygreen}{rgb}{0,0.6,0}
\definecolor{mygray}{rgb}{0.5,0.5,0.5}
\definecolor{mymauve}{rgb}{0.58,0,0.82}
\definecolor{lightgray}{gray}{0.95}

% # <<opts_chunk, include = FALSE>>=
% #   opts_chunk$set(fig.width=10, fig.height=4, out.width = '0.9\\linewidth', echo =FALSE)
% # @

%%%%%%%%%%%%%%%%%%%%%%%%%%%%%%%%%%%%%%%%%%%%%%%%%%%%%%%%%%%%%%%%%%%%%%%%%%%%

\begin{document}


\lhead{University of Geneva\\ \textbf{GLAM}\\
Pr. Eva Cantoni}
\rhead{GSEM\\ Spring 2022 \\ \textbf{Homework 02}}

\begin{center}
\bigskip

\bigskip

\bigskip

\bigskip

\bigskip

\bigskip

\fbox{{\Large Handwritten exercises on GLMs - Handout}}

\bigskip

\bigskip

\bigskip
\end{center}


%%%%%%%%%%%%%%%%%%%%%%%%%%%%%%%%%%%%%%%%%%%%%%%%%%%%%%%%%%%%%%%%%%%%%%%%%%%%

\setlength{\parskip}{1ex plus 0.5ex minus 0.2ex}
\parindent=0cm
\setlength{\parskip}{1ex plus 0.5ex minus 0.2ex}
\linespread{1.1}
\sloppy

\section{\underline{Exercise 1}}
Show that the Binomial distribution and the Gamma distribution belong to the exponential family of distributions, by identifying $\theta$, $A$, $\phi$, and the functions $b()$ and $c()$ (cf.\ p.~53 of the course slides). Using the properties of the distributions of the exponential family, give the expectation and the variance in each case.


\textit{Notation:}
\begin{itemize}
\item If $Y$ follows a  binomial distribution ${\mathcal B}(m,p)$, then 
$$P(Y=y; m,p) = {m \choose y} p^{y} (1- p)^{m-y}, \ \  \mbox{ for } y=0, \ldots, m.$$

{\color{blue} 

From the pmf above, we can write

\begin{align*}
  \mathbb{P}(Y = y ; m, p) &= {m \choose y} \left(\frac{p}{1-p}\right)^{y} (1- p)^{m}\\
&=\exp\{\ln{m \choose y}+y\ln \left(\frac{p}{1-p}\right) + m\ln (1-p)\}\\
&= \exp\{y\ln \left(\frac{p}{1-p}\right) + m\ln (1-p) + \ln{m \choose y}\}  {\color{red}} 
\end{align*}
From this, we deduce $\theta = \ln \frac{p}{1-p}$ {\color{red}} . Also, $b(\theta) = - m \ln(1-p)$.


 From $\theta = \ln \frac{p}{1-p}$ we obtain

\begin{align*}
    \exp \theta&=\frac{p}{1-p}\\
    \exp\theta - p\exp\theta&=p\\
    \exp\theta&=p(1+\exp\theta)\\
    p&=\frac{\exp\theta}{1+\exp\theta} 
\end{align*}
and therefore $b(\theta)= - m  \ln (\frac{1}{1+\exp\theta}) = m \ln(1+\exp\theta)$. \\

At this point, we have :
\begin{align*}
  \mathbb{P}(Y = y ; m, p) &= \exp\{y\theta - b(\theta) + \ln{m \choose y}\},
\end{align*}
which leaves $\phi=1$, $A=1$ and $c(y, \frac{\phi}{A})= \ln{m \choose y}$
\begin{itemize}
  \item $\displaystyle \theta = \ln \frac{p}{1-p}$
  \item $\displaystyle b(\theta) = m\ln \op 1 + \exp \theta \cp$
  \item $\displaystyle \phi  = 1$ 
  \item $\displaystyle A = 1 $ 
  \item $\displaystyle c \op y , \frac{\phi}{A}\cp = \log { m \choose y}$ 
\end{itemize}

This gives us $\mathbb{E}(Y|\theta)=b'(\theta)=m\frac{\exp\theta}{1+\exp\theta}=mp$  and $\text{Var}(Y|\theta)=m\frac{b''(\theta)\phi}{A}=\frac{\exp\theta}{(1+\exp\theta)^2}=mp(1-p)$ 
}


\item If $Y$ follows a Gamma distrbution $\Gamma(\mu,\nu)$, then 
$$f(y;\mu,\nu)=\frac{\nu \exp(-\nu y/\mu)(\nu y/\mu)^{\nu-1}}{\mu\Gamma(\nu)}, \ \  \mbox{ for } y>0.$$
\end{itemize}

{\color{blue} 

From the pdf above, we can write

\begin{align*}
  f(y;\mu,\nu)&=\exp\left(\ln\frac{\nu \exp(-\nu y/\mu)(\nu y/\mu)^{\nu-1}}{\mu\Gamma(\nu)}\right)\\
&= \exp\left(\ln \nu - \nu \frac{y}{\mu} + \nu\ln\nu+\nu\ln\frac{y}{\mu}-\ln\nu-\ln\frac{y}{\mu}-\ln\mu-\ln\Gamma(\nu)\right)\\
&= \exp\left(-\nu( y\frac{1}{\mu} - \ln\frac{1}{\mu}) + \nu\ln\nu+(\nu-1)\ln y - \ln\Gamma(\nu) \right) 
\end{align*}







\begin{itemize}
  \item $\displaystyle \theta = \frac{1}{\mu}$ 
  \item $\displaystyle b(\theta) = \ln(\frac{1}{\mu}) = \ln \theta$ 
  \item $\displaystyle \phi  = \frac{1}{\nu}$ 
  \item $\displaystyle A = -1 $ 
  \item $\displaystyle c\op y , \frac{\phi}{A}\cp = 
  -\frac{A}{\phi} \ln \op - \frac{A}{\phi}\cp 
  + \op - \frac{A}{\phi} -1\cp\ln y - \ln \Gamma \op -\frac{A}{\phi}\cp  $ 
    \end{itemize}
    
    This gives us $\mathbb{E}(Y|\theta)=b'(\theta)= \frac{1}{\theta}=\mu$   and $\text{Var}(Y|\theta)=\frac{b''(\theta)\phi}{A}=\frac{1}{\theta^2\nu}=\frac{\mu^2}{\nu}$ 
}


\section{\underline{Exercise 2} }

The beta distribution is used for situations where the measured variable takes values between 0 and 1, for example proportions.

We let $Y$ follow a beta distribution with parameters $\mu$ and  $\phi$, where $\mu$ is the parameter of interest and $\phi$ is a nuisance parameter. The density function of $Y$ is
$$f(y) = \frac{\Gamma(\phi)}{\Gamma(\phi \mu) + \Gamma(\phi (1-\mu) )} y^{\phi \mu - 1} (1-y)^{\phi (1-\mu) -1},  \ \  \mbox{ for } 0 < y < 1.$$

Does this distribution belong to the exponential family? Justify your answer, and if yes deduce $E(Y)$ and $Var(Y)$.



{\color{blue} 
We write
\begin{align*}
  f(y)&=\exp\ok\ln\Gamma\op\phi\cp - \ln\Gamma\op \phi \mu \cp - \ln \Gamma\op \phi\op1-\mu\cp\cp +\op\phi\mu-1\cp\ln y + \op \phi\op1-\mu\cp\cp \ln \op 1-y\cp \ck\\
  &= \exp \ok \op \phi\mu-1 \cp\ln y + \op \phi\op1-\mu\cp\cp \ln \op 1-y\cp + \ln\Gamma\op\phi\cp - \ln\Gamma\op \phi \mu \cp - \ln \Gamma\op \phi\op1-\mu\cp\cp \ck. 
\end{align*}
We therefore cannot express $f(y)$ as 
\begin{align*}
  f(y)&=\exp\ok A \frac{y\theta - b\op\theta\cp}{\phi} + c\op y, \frac{\phi}{A}\cp\ck.
\end{align*}

In particular, the term $y \theta$ cannot be identified.

With this parametrisation, the beta distribution does not belong to the exponential family.
  
}
\section{\underline{Exercise 3}}
\begin{itemize}
\item Let $f(y;\theta)$ be the probability density function of a random variable $Y$, for a given parameter $\theta$. Without using the exponential family structure, show that for ${\displaystyle U(\theta;Y) = \frac{\partial}{\partial \theta} \log f(Y;\theta)}$

{\color{blue}
  
}

\begin{enumerate}
    \item{$\exE\left[U(\theta;Y)\right]=0$,}
      {\color{blue}
        \begin{align*}
          \exE\oc U(\theta;Y)\cc  
          &= 
          \exE\oc \frac{\partial}{\partial \theta} \log f(Y;\theta)\cc\\  
          &= 
          \exE\oc \frac{1}{f(Y;\theta)} \frac{\partial}{\partial \theta} f(Y;\theta)\cc \\ 
          &= 
          \int \frac{1}{f(y;\theta)} \oc \frac{\partial}{\partial \theta} f(y;\theta) \cc f(y;\theta) \mathrm{d}y \\
          &= 
          \int  \frac{\partial}{\partial \theta} f(y;\theta) \ \mathrm{d}y \\
          &= 
          \frac{\partial}{\partial \theta} \int f(y;\theta)  \  \mathrm{d}y \quad \textrm{if the support is compact} \\
          &=
          \frac{\partial}{\partial \theta} (1) \quad \textrm{given that $f(y;\theta)$ is a density}\\
          \exE\oc U(\theta;Y)\cc  &= 0
        \end{align*}
      }
    \item{${\displaystyle \exE\left[\frac{\partial U(\theta;Y)}{ \partial \theta}\right]= - \exE\oc U^2(\theta;Y)\cc}$.}
      {\color{blue} 
        \begin{align*}
          \exE\oc\frac{\partial U(\theta;Y)}{ \partial \theta}\cc 
          &=  
          \exE\oc\frac{\partial^2}{\partial \theta^2} \log f(Y;\theta)\cc  
          = 
          \exE\ok \frac{\partial}{\partial \theta}\oc \frac{1}{f(Y;\theta)} \frac{\partial}{\partial \theta} f(Y;\theta)\cc \ck \\
          &= 
          \exE\ok  \frac{1}{f(Y;\theta)^2} \oc  
          f(Y;\theta) \frac{\partial^2}{\partial \theta^2} f(Y;\theta) - \op \frac{\partial}{\partial \theta} f(Y;\theta)\cp^2\cc \ck \\
          &= 
          \exE\oc\frac{1}{f(Y;\theta)} \frac{\partial^2}{\partial \theta^2} f(Y;\theta)\cc
          -
          \exE\oc \frac{1}{f(Y;\theta)} \frac{\partial}{\partial \theta} f(Y;\theta) \cc^2 \\
          &= 
          \int \frac{\partial^2}{\partial \theta^2} f(y;\theta) \ \mathrm{d}y 
          -
          \exE\oc \frac{\partial}{\partial \theta} \log f(Y;\theta) \cc^2 \\
          &= 
          \frac{\partial^2}{\partial \theta^2} \int f(y;\theta) \ \mathrm{d}y - 
          \exE\oc U^2(\theta;Y)\cc 
           \\
          \exE\oc\frac{\partial U(\theta;Y)}{ \partial \theta}\cc  &= -\exE\oc U^2(\theta;Y)\cc
   \end{align*}
      } 
\end{enumerate}

\item What condition do we have to impose on the support of $Y$ to prove these two statements?

{\color{blue}
  The support of the random variable $Y$ needs to be compact, so that in particular the integral  and the partial derivative operators can be exchanged. 
}
\end{itemize}

\section{\underline{Exercise 4}}

In this exercise, we consider the problem of Maximum Likelihood estimation of the parameters of a Poisson random variable with a log link, i.e : 
\begin{equation*}
    g\op\mu_i\cp = \mathrm{ln}\ok \exE \oc Y_i \ | \ \bX_i = \bx_i  \cc\ck = \bx_i^T\bbeta
\end{equation*}

\begin{enumerate}
  \item Write the estimating (score) equations for the proposed model.
  {\color{blue}
  
  Start by writing : 

  \begin{align*}
    \mathcal{L} \op \bbeta; \by, \bX\cp &= \prod_{i=1}^n \frac{\exp(y_ix_i^T\beta)\exp(-\exp(x_i^T\beta))}{y_i!} \\
    \ell(\bbeta; \by, \bX)&=\sum\limits_{i =1}^n [y_ix_i^T\beta - \exp(x_i^T\beta)-\ln(y_i!)]
  \end{align*}
  
  The score function is:
  
  \begin{align*}
   \frac{\partial }{\partial\bbeta}\ell \op \bbeta ; \by, \bX \cp  &= \sum_{i=1}^n \oc y_ix_i-\exp(x_i^T\beta)x_i \cc\\
   &= \sum_{i=1}^n \oc y_i-\exp(x_i^T\beta) \cc x_i 
  \end{align*}
  
  }
  \item Show that 
    \begin{equation*}
   \mathcal J(\bbeta; \by) 
   = \sum_{i=1}^n \exp\{\bx_i^T\beta\} 
   \bx_i \bx_i^T.
   \label{eq:J}
  \end{equation*}
    {\color{blue}
From $ \mathcal J(\bbeta; \by)= -\mathbb{E}\op\frac{\partial^2 }{\partial\bbeta\partial\bbeta^T}\ell \op \bbeta ; \by, \bX \cp\cp$ and what we found in a), 

\begin{align*}
    \frac{\partial^2 }{\partial\bbeta\partial\bbeta^T} \ell \op \bbeta ; \by, \bX \cp&=\frac{\partial}{\partial \beta}\sum_{i=1}^n \oc y_i-\exp(x_i^T\beta) \cc x_i \\
    &= -\sum_{i=1}^n \oc \exp(x_i^T\beta) \cc x_i x_i ^T  \\
    \Leftrightarrow \mathcal J(\bbeta; \by) &= \sum_{i=1}^n \oc \exp(x_i^T\beta) \cc x_i x_i ^T 
\end{align*}
}
  
  \item Write the formula for the Newton-Raphson algorithm.  

  {\color{blue}
\begin{align*}
    \bbeta ^ {(k)}= \bbeta^{(k-1)}+ \ok \sum_{i=1}^n \oc \exp(x_i^T\beta) \cc x_i x_i ^T\ck^{-1}\ok \sum_{i=1}^n \oc y_i-\exp(x_i^T\beta) \cc x_i \ck
\end{align*}
}
\end{enumerate}




\end{document}
